% HackApertus OST Track 2A technical report.
% Japanese working draft. The final submission is planned as an English PDF of at most six pages.
% Build: LuaLaTeX + BibTeX (plain). Bibliography keys live in refs.bib.
\documentclass[10pt,a4paper]{ltjsarticle}

\usepackage[margin=24mm]{geometry}
\usepackage{fontspec}
\usepackage{luatexja-fontspec}
\setmainjfont{HaranoAjiMincho}
\setsansjfont{HaranoAjiGothic}
\usepackage{booktabs}
\usepackage{graphicx}
\usepackage[hidelinks]{hyperref}
\usepackage[section]{placeins}

\newcommand{\TBD}[1]{\fbox{\textsf{TBD: #1}}}
\renewcommand{\refname}{参考文献 (References)}

\title{HackApertus OST Track 2A 技術レポート\\
  \large スイス投票冊子に対する多言語NLIと原文根拠の復元\\
  \normalsize (日本語作業稿)}
\author{チーム名 \TBD{チーム名とメンバー}}
\date{改訂 \TBD{最終提出リビジョンと日付}}

\begin{document}
\maketitle

\begin{center}
  \fbox{\parbox{0.92\linewidth}{\small
  これは2026-10-16 12:40 JST時点の日本語作業稿です。最終提出版は英語のPDFで6ページ以内とする予定です。
  内容は公式テンプレート\cite{reporttemplate}の項目を含みます。
  本稿の数値はすべて、開発中に使った公開データでの測定値です。主催者の非公開held-outデータでの成績はわかっていません。}}
\end{center}

%======================================================================
\section{概要 (Summary)}

本プロジェクトはHackApertus OST Track 2A\cite{ostchallenge}に取り組みます。
入力はスイス連邦の投票冊子(PDF)または短い参照文と、独語、仏語、伊語のいずれかで書かれた主張(claim)です。
システムは主張が原文から導けるか(0=entailment)、情報が足りないか(1=neutral)、原文と矛盾するか(2=contradiction)を判定し、判定の根拠となる原文を返します。

モデルはApertus v1.5の70Bと8Bです\cite{apertus,apertus70b,apertus8b}。
設計の中心は、モデルに引用文を書かせないことです。
スクリプトがPDFをページ単位で抽出し、対象の投票案件のページを選び、原文を最大1200文字の区間(source span)に切ってIDを付けます。
モデルは1件につき1回のネイティブのtool callでラベルとspan IDだけを返し、スクリプトがIDから原文とページ番号を復元します。
答えるモデルは入力の構造で決めます。冊子を読むTask Aは70B、参照文だけのTask Bは8Bです。
最後に、モデルが引用したページの隣接ページ、案件の冒頭見開き、主張と語彙が近いページで根拠を5件まで補います。

2026-10-16の朝、旧版(v5)のページ選択が、正解根拠のあるページの36\%をモデルに渡していなかったことがわかりました。
選択を冊子の見出しに基づく方式に直し、根拠の補完を加えた新版(70B親と8B子の2段構成)は、開発用60問(DEV60)、公開60問(PUBLIC60)、確認用60問のすべてでTask Aのmacro-F1と根拠の一致率がv5を上回りました。
その後、新版の420件では70B親の最終回答が8B子の推奨とすべて一致していたことがわかりました。
そこで子をなくし、1回の呼び出しで答える構成(R)を試しました。
Rは3つの問題集すべてでTask AとTask BのF1、根拠の一致率が新版を上回り、トークンは半分以下になりました。
2026-10-16 12:40 JST頃から、Rを本番構成としています。

%======================================================================
\section{課題と評価 (Task and evaluation)}
\label{sec:eval}

Task A(Document)は冊子PDF全体と投票案件名(vote)を受け取り、Task B(Reference)は短い参照文を受け取ります。
入出力はJSONLで、1件につき1行を返します。
公式のSolution API Guide\cite{apiguide}によると、Task Aでentailmentまたはcontradictionと答える場合は、原文言語の逐語引用を物理ページ番号(1始まり)付きで1件以上返します。
1項目は約5,000文字以内で、採点されるのは先頭5件です。
Task Bの根拠は任意で採点されず、返す場合のページは\texttt{null}です。

公式の評価は非公開held-outでの3クラスmacro-F1で、根拠の正確さ、トークン数、ケースあたりの時間も見ます\cite{ostchallenge,apiguide}。
本稿の比較では、主催者の評価スクリプト(根拠の照合はrapidfuzz\cite{rapidfuzz}の部分一致)をローカルで実行し、次の値を使います。
Task AとBのmacro-F1、Task Aで正解根拠のある問題のうち先頭5件の根拠に正解箇所が含まれた割合(以下Hit@5)、1件あたりのトークン数と時間、形式失敗(モデルの回答を受理できず既定回答に退避した件数)です。

データは公開データセット\cite{dataset}のtrain splitだけです。
DEV60、PUBLIC60、確認用60問はどれもここから選んだ60問で、開発中に使ったデータです。確認用60問は30行をTask AとBの両方の形式で使っています。
確認用60問は一度きりの比較のために使わずに残していましたが、比較の後は使用済みとして扱っています。
したがって、本稿の数値は未知データでの性能の推定値ではありません。

%======================================================================
\section{システム構成 (Architecture)}

\begin{figure}[t]
  \centering
  \small
  \setlength{\fboxsep}{6pt}
  \fbox{\parbox{0.86\linewidth}{\centering
  入力JSONL: 投票案件名、主張、PDFまたは参照文\\[2pt]
  $\downarrow$\\[2pt]
  pypdfでページ抽出 → 見出しによる案件ページ選択(Task A)\\[2pt]
  $\downarrow$\\[2pt]
  最大1200文字のsource spanとIDの作成\\[2pt]
  $\downarrow$\\[2pt]
  1つのモデルが1回の\texttt{submit}(ラベルとspan ID)\\
  冊子(Task A)→ Apertus 70B、参照文(Task B)→ Apertus 8B\\
  根拠のない非中立は1回だけやり直し、だめならneutral\\[2pt]
  $\downarrow$\\[2pt]
  IDから原文とページを復元 → 根拠の補完(Task A、最大5件)\\[2pt]
  $\downarrow$\\[2pt]
  出力JSONL: ラベル、原文根拠、トークン数、ケース処理時間
  }}
  \caption{推論パイプライン。モデルはラベルとspan IDだけを返し、引用文とページ番号はスクリプトが原文から復元します。}
  \label{fig:pipeline}
\end{figure}

図\ref{fig:pipeline}に処理の流れを示します。

\paragraph{PDF抽出とページ選択}
pypdf\cite{pypdf}でPDFをページごとにテキスト化します。ページ番号はPDFの物理ページです。
Task Aでは\texttt{document\_selection.py}の\texttt{vote\_pages}が案件のページを選びます。
連邦の冊子は、各見開きの先頭ページに\texttt{Erste Vorlage}、\texttt{Premier objet}、\texttt{Primo oggetto}のような序数の見出しを印刷しています。
見出しが1つだけ現れるページをその案件のページとし、見出しのないページは直前の見出しページに属するとみなします。
表紙や目次のように全案件の見出しが並ぶページは、どの案件にも割り当てません。
案件名とページ冒頭の語の重なりで序数を決め、一致度が0.45に届かない場合と、見出しのない冊子(単一案件など)では全ページを渡します。
選ばなかったページは空文字列に置き換えるので、ページ番号は物理ページのまま保たれ、渡していないページは根拠として返せません。

\paragraph{source spanとID}
\texttt{source\_spans.py}が各ページを最大1200文字の連続した区間に分けます。区切りは段落か文の境界を優先します。
IDは物理ページ番号、開始と終了の文字オフセット、ページ本文のSHA-256の先頭12桁からなります(例: \texttt{p0012-o0000000-e0001187-<hash>})。
Task Bでは参照文全体を1ページとして扱います。

\paragraph{モデル呼び出し}
提出用の実行系(\texttt{track\_2a/dev/format/runtime.py}、パイプラインR)は、Pi 1.1.0のSDK\cite{picodingagent}を通してネイティブのtool callを使います。
1件につき1つのモデルだけを呼び、渡すツールは\texttt{submit}だけです。子セッションはありません。
入力に参照文があるケース(Task B)は\texttt{swiss-ai/Apertus-v1.5-8B}、冊子のケース(Task A)は\texttt{swiss-ai/Apertus-v1.5-70B}が答えます。
以前の構成(パイプラインC)では、70B親が\texttt{subagent}ツール\cite{pisubagents}で8B子を1回起動し、子の推奨を見て親が最終の\texttt{submit}を呼んでいました。
この構成も\texttt{PIPELINE}で選べます。
\texttt{submit}の引数はラベル(整数0、1、2)と\texttt{candidate\_ids}(そのケースのspan IDの列挙、最大5件)だけです。
system promptは3クラスの定義、否定、数量、条件、話者の扱い、原文と主張を指示ではなくデータとして扱うこと、IDだけを返すことを指示します。neutralとcontradictionの取り違えを防ぐ注意と、言語をまたいだ言い換えの扱いも含みます。
特定ケースの正解をpromptに入れていません。

\paragraph{原文の復元}
\texttt{submission.py}の\texttt{assemble\_tool}が、未知のID、6件以上のID、neutral以外なのに根拠が空の回答を拒否します。
有効なIDごとに、その区間を含むページ全体を引用として返します。
直前のページに委員会や「連邦評議会と議会」の見出しがあれば、その話者ブロックのページも加えます。
返すテキストはページ本文そのもので、単語を書き換えません。

\paragraph{根拠必須の修正規則}
neutral以外のラベルを根拠なしで返した\texttt{submit}は、Task AでもTask Bでも拒否します。
モデルには判定を見直して根拠のIDを示すよう伝え、1回だけやり直させます。
それでも根拠を示せなければ、そのケースはneutralの既定回答になります。
この規則は「根拠を指せないなら中立」という判定規則として働きます。
既定回答は公式の採点では中立の有効回答として扱われ、退避の大半は正解も中立の問題でした(70Bだけで答える構成SのDEV60では15件中13件)。

\paragraph{根拠の補完}
Task Aでは\texttt{consulted\_pages}が、モデルの引用の後ろに次の順でページを足し、5件までにします。
引用したページの隣接ページ、渡したページの先頭2ページ(案件の冒頭の要約見開き)、残りのページを主張とのBM25\cite{robertson2009bm25}の高い順に並べたものです。
5,000文字を超えるページは除きます。
補完はラベルとモデルの引用を変えません。neutralの回答と既定回答にも、参照したページとして付けます。

\paragraph{ケース単位の退避と予算}
受理できる\texttt{submit}がないケースや、そのケースの処理が失敗した場合は、そのケースだけをneutralの既定回答にし、理由を診断用のファイルに残します。
要求数の上限は入力件数の5倍、トークン予算はその要求数に1要求あたりの予約量262,400を掛けた値で、入力件数に比例します。
1ケースの上限時間は300秒です。
\texttt{input\_tokens}と\texttt{output\_tokens}はそのケースの全呼び出しのAPIの\texttt{usage}の合計、\texttt{inference\_time\_ms}は文脈の準備を含むケース全体の時間です。

%======================================================================
\section{開発の経緯と診断 (Development history and diagnosis)}

2026-10-14に、system promptの出力指示を\texttt{submit}ツールの引数の形式に合わせた版(v5)を採用しました(\texttt{docs/prompt-eval-log.md})。
10-15の夜には、検証役や複数エージェントの研究\cite{dhuliawala2023cove,huang2024selfcorrect,du2023debate,smit2024mad}を動機に、8Bに事実やspanを選ばせる役を親に足す3つの構成を試しましたが、形式失敗が増え、根拠の提示も改善しなかったので採用しませんでした。
v5はDEV60の5回平均でTask A F1 0.650、Task B F1 0.815、Hit@5 0.263でした。

10-16の朝、モデルを呼ばずに、DEV60とPUBLIC60でTask Aの正解根拠がある39問を調べました。
14問(36\%)で、正解根拠のページをモデルに渡していませんでした。
旧方式は\texttt{im Detail}などの語で案件詳細節の開始を探していましたが、この語が目次ページに誤って一致し、データセットの参照文の出典になりやすい案件の要約見開き(In Kürze、En bref、In breve)が落ちていました。
これとは別に、3問は正解根拠がpypdfの抽出したどのページとも一致しませんでした。

見出しによる方式に置き換え、確認用の30行を除く公開データで、正解根拠がページと一致する501行を検査しました(表\ref{tab:audit})。
個別の正解行に合わせた調整はしていません。

\begin{table}[t]
  \centering
  \caption{ページ選択の検査(モデル呼び出しなし、501行)。}
  \label{tab:audit}
  \small
  \begin{tabular}{lrrr}
    \toprule
    選択方式 & 正解根拠のページを含む & 文脈文字数 平均 & 最大 \\
    \midrule
    旧(目次の正規表現) & 352 / 501 & 38,764 & 236,824 \\
    新(見出し) & 501 / 501 & 33,463 & 162,973 \\
    \bottomrule
  \end{tabular}
\end{table}

ページ選択だけを直したDEV60の1回目では、Hit@5は0.263のままでした。
正解ページを渡しても、モデルが詳細説明や論点の別のページを引用したためです。関連箇所を入力に含めても使われるとは限らないという報告\cite{liu2023lost,hsieh2024ruler}と同じ傾向です。
そこで根拠の補完を加えました。

本日(コミット650de65)、さらに3つの変更を入れました。
同じIDが重複した回答を修正の往復なしに1件にまとめること、公式の規則に合わせてTask Bの根拠を任意にすること、引用ページの隣接ページを補完に加えることです。
保存済みの回答をモデルを呼ばずに採点し直すと、正解根拠を含む件数は149件中89件から98件に増えました。
次節の表の新版の値は、この変更より前の版によるものです。

新版の7回の実行(DEV60を3回、PUBLIC60を1回、確認用60問を3回、計420件)を調べると、70B親の最終回答は420件すべてで8B子の推奨と一致していました。
親の2回の呼び出し(原文を読んで子を呼ぶ1回と、子の答えを写す1回)が、全トークンの63\%、待ち時間の65\%を占めていました。
そこで、1つのモデルが1回の\texttt{submit}だけで答える構成を実行系に加えました。70Bだけで答えるSと、Task Aを70B、Task Bを8Bに振り分けるRです。

8Bだけの1回呼び出しは、DEV60でTask A F1 0.714、Task B F1 0.808から0.883、8,455トークン/件、2.3秒/件でした。
同じ8Bでも、tool callをやめてテキストのJSONで答えさせると、Task A F1は0.361まで落ちました。
8Bではネイティブのtool callが必須で、それでも冊子の問題では70Bに届かないので、Task Aは70Bに残しました。

Task Bで根拠なしの非中立回答を受け入れていた間は、70Bが話題違いの主張を根拠なしで矛盾と答え、Task B F1が0.693まで落ちました。
コミット8d5c7e6で、根拠必須の規則をTask Bにも戻しました。

%======================================================================
\section{結果 (Results)}

表\ref{tab:results}に、v5、新版(70B親と8B子、コミット76a0e6e)、採用した構成Rの比較を示します。
3つの問題集はどれも開発中に使ったものです。確認用60問も、v5と新版を比べた後は使用済みです。
確認用60問では、v5、新版の順に交互に各3回を実行し、すべて終わってから採点しました。
Rの値は各問題集の1回の実測です。
温度0でも、同じ入力で実行ごとに数件の差が出ます。1回どうしの差が0.03程度なら、偶然の範囲と見ています。

\begin{table}[t]
  \centering
  \caption{v5、新版(70B+8B)、Rの比較。値は平均で、実行回数はv5が5、5、3回、新版が3、1、3回、Rが各1回です。既定回答は全実行の合計で、Rは根拠必須の規則による退避を含みます。PUBLIC60のv5の時間は参照した資料に記録がありません。}
  \label{tab:results}
  \footnotesize
  \setlength{\tabcolsep}{3.5pt}
  \begin{tabular}{lrrrrrrrrr}
    \toprule
     & \multicolumn{3}{c}{DEV60} & \multicolumn{3}{c}{PUBLIC60} & \multicolumn{3}{c}{確認用60問} \\
    \cmidrule(lr){2-4}\cmidrule(lr){5-7}\cmidrule(lr){8-10}
    指標 & v5 & 70B+8B & R & v5 & 70B+8B & R & v5 & 70B+8B & R \\
    \midrule
    Task A F1 & 0.650 & 0.705 & 0.875 & 0.701 & 0.869 & 0.967 & 0.734 & 0.797 & 0.804 \\
    Task B F1 & 0.815 & 0.815 & 0.853 & 0.917 & 0.931 & 1.000 & 0.929 & 0.929 & 0.933 \\
    Hit@5 & 0.263 & 0.544 & 0.74 & 0.30 & 0.65 & 0.75 & 0.569 & 0.625 & 0.75 \\
    トークン/件 & 28,232 & 24,361 & 11,004 & 30,382 & 28,964 & 11,000 & 25,980 & 28,843 & 10,956 \\
    時間/件(秒) & 6.3 & 5.8 & 3.0 & 記録なし & 6.5 & 3.0 & 6.7 & 6.8 & 3.1 \\
    既定回答(件) & 0 & 0 & 6 & 7 & 1 & 4 & 3 & 0 & 3 \\
    \bottomrule
  \end{tabular}
\end{table}

\paragraph{v5と新版}
DEV60では3回とも新版のTask A F1がv5を上回り(0.690から0.728)、Task Aの矛盾の再現率は42\%から53\%に上がりました。
確認用60問のv5の既定回答3件は、文脈が送信上限の240,000バイトを超えた1問(248,883バイト)を各回で数えたものです。
v5はこの1問をトークン0として数えているので、両版が処理した59問で比べるとv5が26,420、新版が26,692トークン/件です。
事前に登録した合格基準のうち、新版は「Hit@5がv5より0.10以上高い」(実際は+0.056)と「トークン/件がv5以下」(実際は+11\%)を満たしませんでした。
所有者はこの2つが大会の規則にない基準だったとして、大会の評価項目で比べ直し、新版を採用しました(\texttt{docs/confirmation-protocol-2026-10-16.md})。

\paragraph{新版とR}
Rは3つの問題集すべてで、Task A F1、Task B F1、Hit@5が新版を上回りました。
ただし確認用60問のTask A F1(0.797と0.804)とTask B F1(0.929と0.933)の差は、偶然の範囲です。
トークンは約11,000/件で新版の半分以下、時間は3.0から3.1秒/件です。
表とは別に、後で新版をDEV60で1回動かすとTask A F1 0.733、Task B F1 0.815、Hit@5 0.68、24,404トークン/件、6.1秒/件で、Rはこれも上回ります。
70BだけのSも、Task A F1とHit@5は3つの問題集すべてで新版を上回りましたが、確認用60問のTask B F1は0.865で新版より低くなりました。
SではなくRを選んだのは、既定回答が少なく(SはDEV60、15件、PUBLIC60、10件、確認用60問、8件)、トークンが約8\%少なく、Task Bで「話題違いを矛盾と答える」70Bの傾向を避けられるからです(\texttt{docs/single-call-results-2026-10-16.md})。

%======================================================================
\section{限界と今後 (Limitations and next steps)}

\paragraph{評価の限界}
すべての数値は開発中に使った公開データでの値で、60問、1回から5回の実行です。採用したRは各問題集1回しか測っていません。
公式のheld-outでの成績はわかりません。

\paragraph{抽出の限界}
前述の新版の420件のうち、Task Aで根拠が採点される149件を調べると、26件は正解の根拠文がpypdfの抽出テキストと一致しません。
この26件は、ページ選択や補完をどう変えても根拠として当てられません。
テキストを抽出できないPDFはOCRを行わずエラーにします。
見出しのない冊子では全ページを渡すので、入力が長くなります。

\paragraph{親モデルの役割}
新版の420件すべてで、70B親の最終回答は8B子の推奨と一致し、親はトークンと時間を増やすだけでした。
この問題は、子をなくした1回の呼び出し(R)で解消しました。

\paragraph{残る分類誤り}
Rの3問題集(各1回、計180問)では、Task Aの誤りは90問中11件で、矛盾を中立とした5件、矛盾を含意とした3件、含意を矛盾とした2件、中立を矛盾とした1件です。Task Bの誤りは90問中6件で、矛盾を中立2件、矛盾を含意2件、含意を矛盾2件です。
根拠規則による退避13件のうち11件は正解が中立で、矛盾を取りこぼしたのは2件でした。
矛盾の取りこぼしが最も多く残っていますが、件数が少ないため、指示文を変える根拠はまだ弱いと考えています。

\paragraph{根拠の性質}
source IDは引用が原文のどこから来たかを保証しますが、引用が判定を支えるかどうかは保証しません\cite{gao2023alce}。
補完で足したページは検索で選んだもので、モデルの判断ではありません。

\paragraph{今後}
抽出が一致しない26件に対して、pypdf以外のPDF抽出ライブラリを比べる予定です。
また、Rを各問題集で複数回実行し、実行ごとのばらつきを見積もります。

%======================================================================
\section{再現方法 (Reproducibility)}

提出物はDockerで動かします。
公式の契約\cite{apiguide}では、CPUのみのlinux/amd64イメージで、注入された\texttt{BASE\_URL}と\texttt{API\_KEY}をすべてのモデル呼び出しに使い、実行時にダウンロードしないことが求められます。
\texttt{Dockerfile.agent}は、Node 26.10.0の段階で\texttt{package-lock.json}からPi 1.1.0と\texttt{@tintinweb/pi-subagents}を入れ、\texttt{python:3.11-slim-bookworm}に\texttt{track\_2a/src}と実行系をコピーします。
入口は\texttt{runtime.py}で、既定は\texttt{PIPELINE=R}です(\texttt{Dockerfile.agent}と\texttt{Makefile}の両方)。
Rではケースごとのモデルを入力の構造で決めるので、\texttt{LLM\_MODEL}を指定しても使いません。
\begin{verbatim}
make run INPUT_DIR=/path/to/input INPUT=/data/input.jsonl \
  OUTPUT_DIR=/path/to/new/output
\end{verbatim}
以前の70B親と8B子の構成は\texttt{PIPELINE=C}、70Bだけの1回呼び出しは\texttt{PIPELINE=S}で選べます。
\texttt{BASE\_URL}と\texttt{API\_KEY}は主催者が注入する値を使い、ホストの鍵は読みません。
出力ディレクトリには正規のJSONLのほか、退避の理由を残す\texttt{.diagnostics.jsonl}と、評価用の\texttt{.evaluation.jsonl}を書きます。
\texttt{track\_2a/dev/format/README.md}によると、開発ホストではDocker daemonを使えず、イメージのbuildとrunは確認できていません。

開発用のハーネスでは、Pi 1.1.0を\texttt{\textasciitilde/.cache/ost-pi-1.1.0}に置き、\texttt{PI\_PACKAGE\_ROOT}と\texttt{PI\_SUBAGENTS\_ROOT}でその場所を指定します。
リポジトリURLと提出リビジョンは\TBD{リポジトリURLと最終コミット}、デモは\TBD{デモ}です。

%======================================================================
\section*{ライセンス (License)}
\addcontentsline{toc}{section}{ライセンス (License)}

リポジトリのコードにはApache License 2.0が付いています(リポジトリ直下の\texttt{LICENSE})。本レポートのライセンスはCC-BY-4.0とします。
大会規約では、コードと重みはApache-2.0、文書はCC-BY-4.0、提出データセットはCDLA-Permissive-2.0での許諾が定められています。
同梱したpypdfは、同梱のライセンスファイルの条件に従います。
公開データセット\cite{dataset}のカードにはMITと記載され、作成者はOSTのAI Applications and Deployment Labです。投票冊子PDFの再配布権は別途確認が必要で、本レポートには同梱しません。

%======================================================================
\begingroup
\small
\bibliographystyle{plain}
\bibliography{refs}
\endgroup

\end{document}
